\documentclass[12pt]{article}

\usepackage[utf8]{inputenc}
\usepackage{amsmath, amssymb, amsthm}
\usepackage{geometry}
\geometry{a4paper, margin=2.5cm}

\title{Sigma-Algebră}
\author{}
\date{}

\theoremstyle{definition}
\newtheorem{definition}{Definiție}

\begin{document}

\maketitle

\begin{definition}
O familie de mulțimi $\mathcal{F}$ se numește \emph{sigma-algebră} dacă:
\begin{enumerate}
    \item $\Omega \in \mathcal{F}$.
    \item Dacă $A \in \mathcal{F}$, atunci $A^{c} \in \mathcal{F}$.
    \item Dacă $(A_n)_{n \ge 1}$ este o familie de mulțimi din $\mathcal{F}$, atunci
    

\[
        \bigcup_{n=1}^{\infty} A_n \in \mathcal{F}.
    \]


\end{enumerate}
\end{definition}

\noindent
Sigma-algebre sunt închise la reuniuni numărabile.

\end{document}
